
% language=us runpath=texruns:manuals/ontarget

\startcomponent ontarget-signing

\environment ontarget-style

\startchapter[title={Signing documents}]

\subject{How did we end up here}

Here I discuss a feature of \PDF\ documents that \TEX\ users can safely ignore
most of the time but that is part of the package. But before we arrive at
describing what it is, first a few words about \PDF\ and the way it is used.

When \PDF\ showed up as a format, \DVI\ was what \TEX\ users had to deal with. It
was possible to preview the result with a \DVI\ viewer or convert it to some
format suitable for a printer. The \DVI\ format is rather minimal and has no
resources like fonts and images embedded. Basically the file only positions
glyphs and rules on a canvas and the glyphs are references to external fonts.
Color and similar effects have to be implemented using the \type {\special}
primitive that puts directives for the backend driver in the file. Although
technically one could embed fonts and images using specials and thereby create a
self|-|contained file it never happened, partly because it demands a dedicated
viewer and (definitely at that time) increased runtime.

At that time \POSTSCRIPT\ was one of the popular output formats. For a long time
I used \DVIPSONE\ for (robust) printer output (with outline fonts) and \DVIWINDO\
for previewing (because it was fast, supported color, graphics and hyperlinks).
The initial road to \PDF\ was to add another step to the conversion: using
\ACROBAT\ to convert a \POSTSCRIPT\ (enhanced with so|-|called pdfmarks) into a
\PDF\ file. Later direct \DVI\ to \PDF\ converters showed up alongside \PDFTEX\
that integrated an alternative backend. We can realize that without these more
direct methods \TEX\ would not be as popular as it is now. The fact that in
\CONTEXT\ we had a rather generic (abstract) backend, where specific drivers
plugged in, indicates that we didn't foresee that \PDF\ would eventually take
over.

One of the reasons \PDF\ showed up alongside \POSTSCRIPT\ is that it removes the
interpretation part from the end result. Where \POSTSCRIPT\ is a programming
language and the result needs to be interpreted in the printer or in a viewer
(like \GHOSTVIEW), \PDF\ is a collection of related objects that express what
gets rendered in what way. Often the \PDF\ files are smaller, also due to
compression (so they transfer faster), and the lack of additional processing
removes a bottleneck in high speed and high resolution printing. If you look at
it this way, \PDF\ is primarily a {\em printer format}.

Some tools in the Adobe suite of editing programs use(d) a mix of binary and
\POSTSCRIPT\ to store the state. At some point \PDF\ became the container format,
although one could find curious mixes of \POSTSCRIPT, \PDF, even configurations
stored as typeset streams, but that might have been normalized by now. So, from
this perspective \PDF\ is a {\em storage format}, kind of an object|-|oriented
one.

When \PDF\ showed up it had some rudimentary support for annotations, like
hyperlinks, and also for embedding of audio and video. I'm pretty sure that \TEX\
was among the first programs to support this. Over time more annotation types
showed up, like widgets (forms), complex media and 3D, \JAVASCRIPT, comments and
attachments. Widgets evolved a bit and one has to play rather safe (and not use
all features) because some bugs and side effects became features and there is
limited support in viewers, if only because often \JAVASCRIPT\ is assumed. Media
have always been sort of a mess, especially when the straightforward annotations
were dropped. We have always supported them but I consider them unreliable in the
long run. Most hyperlinks are working as expected, comments (a form of \PDF\
annotations) when used wisely also work ok, although that depends on the viewer
(interfaces change). No matter what one thinks of all this, here \PDF\ is
definitely a {\em viewing format}. Because comments can be added, \PDF\ files can
also be used in redacting workflows.

It is also worth remembering that in the early days only Acrobat was available
for viewing; there was even a version for \MSDOS. The Reader only offered basic
functionality and for the real deal one needed Exchange, which didn't come free.
There was a complicated scheme for shipping a special version with documents:
basically that was the business model for using \PDF\ for an on|-|screen reading
experience. This was not a success (cumbersome as well as expensive), and when
Internet access became more common, using \CDROM\ for distributing documents was
soon obsolete, let alone installing a related closed source and operating system
dependent viewer. We'll see that with document signing, another attempt at a
business model is made.

Following up on that we now arrive at two additional aspects. One is
accessibility and I could spend a whole article on that. Let's stick to the
observations that the idea is that rendered content can be reinterpreted for
reading aloud, for reflow (sic), maybe for cut|-|and|-|paste. Personally I wonder
why one would use \PDF\ to provide adaptive accessibility, because \HTML\ is
meant for that. Distributing a structured source might be more efficient when
interpretation is needed. Anyway, it's not that hard to support but we end up
with a bloated \PDF\ file, while the \PDF\ format started out with attempts to
minimize size. I understand that publishers don't like to distribute sources but
if this is the solution one can wonder. The second addition is to render and
present text in a way that cannot be tampered with and this is where signing
comes in: can we somehow mark a document such that the receiver can trust its
content. More about that later, but what we can say here is that \PDF\ has become
a {\em distribution format}. Conforming to standards, tagging, encryption and
signing all play a role in this.

No matter what usage we consider, all of them depend on reliability. Can we show
and process a document in the future? Can \PDF\ be relied upon as a long|-|term
archival format? From that perspective, standardization has to be mentioned.
Already early in the history of \PDF, plugins (and additional workflows) provided
the printing industry ways to check if a file was okay. One should think of color
spaces being used, fonts being present, etc. Some tools manipulated the \PDF, not
always with the best outcome, but we leave that aside. Other tools were a bit
more tolerant than might be considered healthy, for instance by ignoring a bad
xref table (basically the registry of objects in a \PDF\ file) and either fixing
it or just generating one from scratch. Although Acrobat can complain or fail to
open a document, on the average commercial and open source tools are tolerant
enough and the lack of a proper error log means that the \PDF\ generators don't
get fixed when that still can be done. It is also good to keep in mind that there
is a whole industry around \PDF\ generation, validation, manipulation, etc.\ and
huge money making machines are not always on the retina of, for instance, \TEX\
users who produce \PDF. Of course by now there are so many documents in \PDF\
format around that being tolerant kind of comes with the package. Validators like
\VERAPDF\ evolve and a document that is ok today (2023) might fail the test
tomorrow, and the verdict even depends on the \PDF\ framework being used (there
are options). Where \TEX\ users can often regenerate a document from source this
is not true for the majority of documents produced elsewhere.

It is also important to notice that rather soon in the history of \PDF,
\GHOSTSCRIPT\ became an option for viewing and at some point commercial and open
source viewers showed up. Not all were perfect and even today there are
differences in quality and functionality. A good test is how well
cut|-|and|-|paste deals with spaces and how well a test area gets selected. The
open source viewers are slow in catching up, but because the evolution of media
\PDF\ annotations isn't that stable either for most purposes viewers like
\SUMATRAPDF\ (\MSWINDOWS) or Okular (\LINUX) is what I use today, especially now
that Acrobat has moved to the cloud. There is also some competition from browsers
that show \PDF. For purposes of signing (which we'll get to next) one probably
has to rely on Acrobat for a while, but we'll see.

So, what does signing bring to this? Digital signatures have been around for a
while. You can for instance sign a document with a certificate (similar to what
secure webservers do with sites). In that case the distributed blob has
security|-|related information as well as the content. A validating application
can take the content and check if it has been tampered with. It can do so off
line (with limited security) but also go online and check the embedded
certificate. With signed \PDF\ files the same is true apart from the fact that
here the signature is a partial one, not embedding the data. Instead the
signature is embedded in the \PDF\ file.

Before we move on we have to stress that signing is not the same as (password
protected) encryption. A signed \PDF\ file is by default just readable, unless
one explicitly encrypts the file. These processes are independent. Here we ignore
encryption; suffice it to say that \CONTEXT\ can do it, but apart from users
asking for it I don't know if it ever gets applied. We discuss the process of
signing in the perspective of \CONTEXT, although in itself it is not bound to
that macro package.

Let's first look at how text ends up in a \PDF\ file. Take this source file, in
\CONTEXT|-|speak:

\starttyping
\startTEXpage[offset=1dk]
   some text
\stopTEXpage
\stoptyping

This leads to a so|-|called page stream that contains this (except normally you
are likely to see garbage because compression is applied, so decompress first):

\starttyping
BT
/F1 10 Tf
1.195517 0 0 1.195517 7.485099 7.616534 Tm
[<000100020003000400050006000400070006>] TJ
ET
\stoptyping

We switch to a font (\type {Tf}) with id \type {F1}, set up a text transform
matrix (\type {Tm}) and render the four plus four characters (\type {TJ})
indicated by their index into a (in our case subsetted) font. The \type {0005} is
not really a character: it refers to a space. It looks unreadable but one can
figure out the text by consulting the \type {ToUnicode} resource associated with
the font as it has the mapping from the index numbers to \UNICODE\ (with comments
added):

\starttyping
<0001> <0073> % s
<0002> <006F> % o
<0003> <006D> % m
<0004> <0065> % e
<0005> <0020> %
<0006> <0074> % t
<0007> <0078> % x
\stoptyping

There is nothing hidden here and one can actually even change the text by
changing an index in the page stream, although of course you can only use indices
that are available and you also have to accept weird rendering due to the change
in progression when the referenced glyph has different dimensions. More extensive
tampering with the document has more severe consequences. For instance, the page
object looks like this:

\starttyping
3 0 obj
<< /Length 118 >>
stream
...
endstream
endobj
\stoptyping

So changing the content also demands changing the \type {Length}. Even worse,
there is an entry in the object cross reference table:

\starttyping
0000000075 00000 n
0000000244 00000 n
\stoptyping

that needs to be adapted, including all following entries. So, tampering is
possible but not something that is likely to happen. Nevertheless we continue as
if some guard against this is needed. We now assume the following document:

\startbuffer[nocompression]
    \nopdfcompression
\stopbuffer

\startbuffer[signed]
\setupinteraction[state=start]

\definefield[signature][signed]

\defineoverlay[signature][my signature]

\starttext
    \startTEXpage[offset=1ts,frame=on,framecolor=darkblue]
        sign: \inframed
          [background=signature,framecolor=darkred]
          {\fieldbody[signature][width=3cm,option=hidden]}
    \stopTEXpage
\stoptext
\stopbuffer

\typebuffer[signed]

We get a small, one page, document:

\startlocallinecorrection
    \typesetbuffer[nocompression,signed]
\stoplocallinecorrection

This document has a widget that looks like this (some less relevant entries are
omitted):

\starttyping
2 0 obj
<<
    /Type    /Annot
    /Subtype /Widget
    /FT      /Sig
    /T       <feff007300690067006e00610074007500720065>
    /V       1 0 R
    /Rect    [ 38.445125 11.522571 123.484489 23.471172 ]
>>
endobj
\stoptyping

In principle we could add an appearance stream and decorate the widget but when
adding signature support to \CONTEXT\ I found that using a parent|-|kid approach,
for instance, was not appreciated by some programs (I used mutool (mupdf), pdfsig
(poppler), Okular and Acrobat Reader for some basic testing), so in the end the
\type {V} key ended up in the root widget. It probably relates to fuzzy
specifications, experiments with specific tool chains, non|-|public validation
processes, etc. Round trip signing and verification seems not entirely trivial,
so best to play safe.

When the value of a \type {Sig} widget is a string, signing is up to the viewer
but when we have a dictionary the signature can be in the file. The \type {V}
value of \type {1 0 R} is a reference to a dictionary with object number~\type
{1}. Here is what that value looks like when we generate this document:

\starttyping
1 0 obj
<<
    /ByteRange [ 2000000000 2000000000 2000000000 2000000000 ]
    /Contents  <00000000000000000000....00000000000000000000>
    /Filter    /Adobe.PPKLite
    /SubFilter /adbe.pkcs7.detached
    /Type /Sig
    >>
endobj
\stoptyping

The \type {Filter} and \type {SubFilter} entries are sort of default, though
alternatives are possible, which then requires additional information to be added
and also a viewer (or validator) able to deal with it. We leave that aside. The
\type {Contents} hex|-|encoded string is a placeholder for the signature and in
our case is 4096 bytes long. We could compute a bogus signature and check the
size instead. Here the \type {ByteRange} and \type {Contents} are actually
invalid but viewers are (supposed to be) tolerant so it triggers no error. After
all this widget is only consulted when signatures are checked. The general
structure of the file is like this:

\starttyping
%PDF-1.7
....
1 0 obj
<< /ByteRange [ .... ] /Contents <....> .... >>
endobj
....
xref
\stoptyping

The signature ends up between \type {<} and \type{>} and has to be calculated
over the bytes specified by the \type {ByteRange} entries. Although one might
think that these can be arbitrary, in practice it looks like it is best sticking
to the recommendation:

\starttyping
start of file
length upto position < of contents
position after > of contents
length upto end of file
\stoptyping

So basically all except the \type {Contents} value is taken into account. Because
the ranges are part of that, they need to be filled in properly, something that
has to be done after the \PDF\ file is finished because only then is the size
known. In \CONTEXT\ that could be done as part of the main run, but it makes
little sense because we need to adapt the file anyway. If the file is called
\type {sign-001.tex}, we get this:

\starttyping
mtxrun --script pdf --sign --certificate=sign-001.pem --password=test sign-001
\stoptyping

The script will set the byte ranges and fill in the content. It does that by
making a data file and running \type {openssl} with the appropriate parameters,
although with \type {--library} one can avoid the temporary file and gain a bit.
Just for the record: we don't depend on that library but have only a minimal
delayed binding to a few functions, with \LUA\ wrappers so it has no impact on
(compiling) the \LUAMETATEX\ binary. Eventually one ends up with something like
this (values abridged):

\starttyping
1 0 obj
<<
    /ByteRange [ 0000000000 0000006276 0000010375 0000000380 ]
    /Contents  <3082061a06092a864886....010702a082060b308206>
    ....
endobj
\stoptyping

Verifying can be done as follows:

\starttyping
mtxrun --script pdf --verify --certificate=sign-001.pem --password=test sign-001
\stoptyping

which reports:

\starttyping
sign pdf | signature in file 'sign-001.pdf' matches the content
\stoptyping

while changing a byte in the trailer id results in:

\starttyping
sign pdf | signature in file 'sign-001.pdf' doesn't match the content
\stoptyping

For verifying we can load the \PDF\ file and use the \type {ByteRange}
specification but for signing this is less trivial: when we load a \PDF\ file we
load a structure that is ignorant of the position in the file. We could use the
cross reference table to find the position in the file of the object but that
assumes that this table is available. So here we have two alternatives. We can
write an auxiliary file (\type {sign-001.sig}) at the end of the \TEX\ run that
has the relevant information. This approach permits us to keep the \PDF\ file
simple: we reserve enough characters for the ranges and content so we can
overwrite them. If the file is lacking, the sign routine tries to locate the
object in the \PDF\ from the list of widgets and once we know its number we also
know where the object is in the file. This alternative adds a little overhead
because at least the cross reference table has to be loaded. Whatever route we
take, it is still prettier than appending additional objects to the file and
basically creating a new version, which not only makes the file larger but also
keeps unused objects around. Applications like mutool and Acrobat prefer that
route, though, in part because they add their own appearance streams.

We now need to discuss these certificates and that is where it becomes less
convenient. For testing, I use a Let's Encrypt certificate but these officially
cannot be used as they are flagged as web certificates. There is (what's new
here) a whole industry behind this signing. You need to get a certificate
someplace and for that often have to sign up for a yearly subscription. In the
worst case you get a token instead of a file and then have to set up some
delegated workflow. Feeding a document into a \USB\ token is not the most
efficient of all processes, so you will find alternative solutions where you end
up with a dedicated machine in a server rack. This all makes it a no|-|go for a
low or zero budget situation. It also means that for just {\em printing}, {\em
viewing} or {\em storing} purposes signing doesn't make that much sense: it only
adds overhead.

One can argue that signing is not that robust anyway. Just like we can add a
signature to a file, so can anybody. It's all about trust. When a byte in a \PDF\
file is changed validation fails anyway so that is already a signal; we don't
need to verify the certificate for that. And it's not that hard to let a user
upload a file to the origin and let it validate there where the private key is
known. But wait, isn't it more convenient to do that without uploading? Sure, but
here are some pitfalls. First of all, who knows if a certificate is still valid?
An organization has to spend quite some money on it yearly. And (even root)
certificates expire so in the end the document refers to something invalid
anyway, which effectively makes the document expire after some time. Saving
documents and providing them again might be cheaper and also has the advantage of
archiving. For long term archiving signing makes little sense anyway (expiration,
cracking).

So why do we bother to add signing to \CONTEXT ? The answer is simple: user
demand. Just like being forced to use some \PDF\ standard, users can be forced to
comply with what the organization prescribes with respect to signing, even when
in the end it's just a demand, and nothing is actually done with it. So the main
question is: after showing that it can be done, what eventually happens; how does
the workflow look? It's comparable to tagging: it is sometimes demanded, but
after that the lack of useful tools make it just a box to be ticked.

There are other alternatives to making users feel good about a document: provide
a printed copy, keep the original someplace for downloading, maybe make it
possible to regenerate a document from source, maybe even provide the resource.
Generate a string hash and keep that available alongside the original. In the end
it is all about trust, indeed.

Let's end on a positive note. Getting to know what has to end up in the file is
not that trivial and as with much on the Internet, looking for solutions quickly
brings you to a subset of partial and sometimes confusing answers and solutions.
This is why in the end I decided to just look in the code base of \type {openssl}
that comes with examples and eventually one can sort out something not too
complex. One of the interesting observations was that the binary blob is a
structured key|/|value sequence using technology from decades ago (1984), when
data had to be transferred reliably between architectures and programming
languages: Abstract Syntax Notation One (\cap{ASN}.1). It makes old \TEX ies
feel young when old tools survive beyond the modern short lifespan of fancy web
technologies. I might eventually spend some more time on this, just for the fun
of it.

If you want to know more details: the official \ISO\ standard on \PDF\ has some
sections on the matter; a more comprehensive summary can be found in \quotation
{Digital Signatures in a \PDF, Adobe Systems Incorporated, May 2012}. There is
also \quotation {\cap{CDS} Certificate Policy, Adobe Systems Incorporated,
October 2005} but I suggest to ignore that one unless you're forced to implement
the more expensive route.

Some final words on the mentioned formats. For printing and storage this feature
is not needed. Nor for regular viewing, because users probably don't care that
much if a manual or book is signed and it's unlikely that certificates last that
long (or stay secure for that matter). But it might make sense for distributing
documents with some legal meaning in the short term. In that perspective having
this feature in \CONTEXT\ makes most sense in specific workflows. But it doesn't
hurt to know that \TEX\ is still able to adapt itself to these situations.

{\em Many thanks to Karl Berry for copy|-|editing this text!}

\stopchapter

\stopcomponent

% context sign-001
% mtxrun --script pdf --sign   --certificate=sign-001.pem --password=test --uselibrary sign-001
% mtxrun --script pdf --verify --certificate=sign-001.pem --password=test --uselibrary sign-001
